% LSPV-5 SYNDROMEGUARD1, 2026-10-15 12:54 KST
% Original proposal card sha256 85179325e2be03dece73b934207d30c7efb9edfbf33326544fc834da9768e496
\documentclass[11pt]{article}
\usepackage{amsmath,amssymb,amsthm,url}
\newtheorem{theorem}{Theorem}
\newtheorem{proposition}[theorem]{Proposition}
\title{SYNDROMEGUARD1: Source-Checked Conservation Certificates Across Exponentially Large Configurations}
\author{LSPV-5}
\date{15 October 2026, 12:54 Asia/Seoul}
\begin{document}\maketitle
\noindent\textbf{Status: PROVED/EXACT for a restrictive, protected-memory ordinary-machine source language, independent review pending; COMPUTER TESTS: six independent checkpoints, 6,393 cases and 34,502 assertions, zero detected discrepancies. UNIVERSAL PFCERT AND P vs NP: OPEN.}
The full 26-kB expanded source, executable audit and raw JSON are maintained as separate local artifacts; the present remote edition is self-contained with exact proofs and limitations.

\section{IDEAS PROPOSED BEFORE LITERATURE}
The original candidate register was written before new external literature searching and frozen at SHA-256 85179325e2be03dece73b934207d30c7efb9edfbf33326544fc834da9768e496; these proposals are kept separate from revisions below.

\paragraph{A, CANDIDATE.} Define a polynomial-description syndrome over a protected memory of $2^d$ bits and permit nonlinear data-dependent local XOR toggles only when their aggregate syndrome is zero. Missing hypothesis: every protected write, including generated instructions and address collisions, is mediated. Falsifier: an unguarded STORE, COPY, or alias.

\paragraph{B, CANDIDATE.} Enforce a single source-checked public output $G(\sigma(X),e)$ with one unconstrained bit $e$. Canonically derive both candidate outputs from initialization, not arbitrary user-described proof envelopes. Missing hypothesis: no unmediated output channel. Falsifier: three outputs from different proofs.

\paragraph{C, CANDIDATE.} Compile the guarded source language to a real deterministic multitape machine whose configurations may contain all $2^d$ protected cells, but whose short invariant certificate never expands them. Missing hypothesis: checked interpreter and all program/verification costs kept separate. Falsifier: verifier requiring enumeration of all cells.

\paragraph{D, CANDIDATE.} Recheck matched-oracle Hirahara $K^t$ immunity and collapsing BGS oracle against universally covering two-output PFCERT. Missing hypothesis: same oracle for universal machine, verifier and deciders, uniformly polynomial proofs. This is recovered prior research, not a novel theorem.

\section{LITERATURE CHECK AND REVISIONS}
Hirahara's original ECCC TR20-050, Theorem 3.9, proves P-immunity of the superpolynomial time-bounded Kolmogorov random-string language; its Theorem 3.10 concerns only P-uniform ACC$^0$ immunity for the distinct Levin $Kt$ measure. Original Baker--Gill--Solovay (SIAM J. Comput. 1975) uses oracle query machines and constructs a recursive collapsing oracle $O$ with $P^O=NP^O$. Karr's affine relationships, Murata's Petri-net structural invariants, and Necula's proof-carrying code are antecedents; the transfer to XOR-protected memory below is established directly and no literature novelty claim is made. Aaronson--Wigderson algebrization remains separate.

A survives only with a physically enforced COMMIT guard and XOR alias aggregation. B survives only when terminal output cannot bypass the canonical checked channel. C survives as a real large-memory restricted ordinary-machine construction, with no fast-time guarantee. D remains a previously established project oracle-failure diagnostic and does not turn A--C into a nonrelativizing mechanism.

\section{CONSTRUCTION OR PROOF ATTEMPT}
Fix a source $S$ of $N$ explicit bits and $1\le d,r,n\le N$. A protected array $X\in\mathbb F_2^{2^d}$ is addressed by $a\in\{0,1\}^d$. Immutable circuits $\chi:\{0,1\}^d\to\mathbb F_2^r$ and $G:\mathbb F_2^{r+1}\to\mathbb F_2^n$ have polynomial total source size. Initialization toggles a listed multiset $I$ of at most $N$ addresses, leaving unlisted bits zero. All other private computation, branching, readings, interpretation and generated bytecode can be nonlinear and arbitrarily long. Only a trusted finite-instruction virtual interpreter can change the protected array; no raw protected STORE or unrestricted output primitive exists. Its protected COMMIT accepts $(a_j,\delta_j)_{j=1}^{c}$ with $c\le N$ and computes
\[
\Delta=\bigoplus_{j=1}^{c}\delta_j\chi(a_j).
\]
If $\Delta\ne0$ there is no protected mutation or public output. Otherwise it performs the XOR toggles $X_{a_j}\leftarrow X_{a_j}\oplus\delta_j$ with coincident addresses combined mod 2. The only public HALT output scans the full protected memory to compute
\[
\sigma(X)=\bigoplus_{a\in\{0,1\}^{d}}X_a\chi(a)
\]
and emits $G(\sigma(X),e)$, where $e$ is one arbitrarily computed private bit. That scan may take exponential time. This is an ordinary deterministic source class, not a claim about arbitrary U programs or fast random-access RAMs.

\begin{theorem}[Exact conservation]
Every finite legal execution prefix satisfies
\[
\sigma(X_t)=q_0:=\bigoplus_{a\in I}\chi(a).
\]
This holds for all nonlinear controls, repeated reads, arbitrarily large configurations and aliased commit addresses.
\end{theorem}
\begin{proof}
Initialization is exactly the prescribed XOR toggles. Private steps and rejected commits leave $X$ unchanged. For a successful commit define $\eta_a=\bigoplus_{j:a_j=a}\delta_j$. Then
\[
\sigma(X')\oplus\sigma(X)=\bigoplus_a\eta_a\chi(a)
=\bigoplus_j\delta_j\chi(a_j)=\Delta=0.
\]
Induction proves the claim without enumerating $2^d$ cells.
\end{proof}

\begin{theorem}[Canonical two-output verifier]
There is one ordinary deterministic multitape verifier with conservative simultaneous bounds
\[
T_V(N)=O(N^6),\qquad S_V(N)=O(N^3),\qquad |\pi|=O(N),
\]
whose complete accepted output projection for each valid checked source $S$ is exactly
\[
\mathcal Y_S=\{G(q_0,0),G(q_0,1)\},\qquad
|\{y:\exists\pi\ V(S,y,\pi)=1\}|\le2.
\]
Every physical terminal output of the source belongs to $\mathcal Y_S$.
\end{theorem}
\begin{proof}
The verifier parses only the explicit immutable checked-VM source, the list $I$, the two circuits and the candidate output. It evaluates $\chi$ on each listed address to reconstruct $q_0$, then evaluates $G$ twice and checks whether the candidate equals either canonical string. Optional supplied certificates are rejected unless they agree with the independently calculated $q_0$; an empty certificate suffices. Thus every accepted proof refers to the same list, and the all-proof global multiplicity bound follows. The conservation theorem and checked sole output channel prove terminal coverage. All verifier operations use explicitly represented polynomial-size objects; $X$ is never expanded. The output machine itself may have large time and space.
\end{proof}

\begin{theorem}[Genuine exponential configuration]
For each $d\ge2$, a checked program with $O(d+r+n)$ source bits can physically write all $2^d$ protected cells and terminate with all bits one while preserving a fixed $r$-bit syndrome for any $1\le r\le d$.
\end{theorem}
\begin{proof}
Choose $\chi(a)=(1,a_0,\ldots,a_{r-2})$, excluding the highest address bit. In a $d$-bit counter loop, for every $0\le j<2^{d-1}$ execute the guarded two-toggle COMMIT on addresses $j$ and $j+2^{d-1}$. These have equal $\chi$, hence zero combined syndrome. They are disjoint pairs covering all cells. Starting at zero, each final cell is one. The source is short, but the actual computation and configuration are exponential.
\end{proof}

\begin{proposition}[Exact coNP consecution barrier]
Given arbitrary Boolean circuit $C(z)$, determining whether the unguarded one-step update $x'=x\oplus C(z)$ preserves $q=x$ for every $(x,z)$ is coNP-complete.
\end{proposition}
\begin{proof}
A violating $(x,z)$ is an NP-checkable counterexample. Invariance for all inputs holds exactly when $C(z)=0$ for all $z$, namely Circuit-UNSAT, giving coNP-hardness and membership.
\end{proof}
The checked-VM runtime guard avoids this unrestricted source-verification demand by enforcing safety during execution; it does not extract free invariants from arbitrary source code. A raw assignment $X_0\leftarrow1$ to an initially zero cell falsifies the invariant for $\chi\equiv1$; two SETs at one address do not cancel as two XOR toggles do.

\paragraph{Matched-oracle diagnostic (recovered earlier theorem).}
For one efficiently universal $U^O$, fixed oracle $O$, and $t(n)=n^{\omega(1)}$ dominating fixed polynomial simulation overhead, let $R^{t,O}=\{x:K_U^{t,O}(x)\ge |x|-1\}$. The usual prefix-enumeration argument relative to the SAME $O$ proves this language $P^O$-immune: an infinite $P^O$ subset gives a $n-O(\log n)$-bit prefix, plus $O(\log n)$ length advice, from which polynomially many suffix trials recover a member in time below $t(n)$ and description shorter than $n-1$. At a collapsing BGS oracle $O$ with $P^O=NP^O$, if one uniformly polynomial-length and polynomially oracle-verifiable two-output schema covered every $t$-short $n$-bit endpoint of every $|p|\le n-2$ program, then $E_n=\{y:\exists p,\pi\,V^O(n,p,y,\pi)\}$ would be in $P^O$ and contain at most $2(2^{n-1}-1)=2^n-2$ outputs. Its infinite $P^O$ complement is contained in $R^{t,O}$, contradiction. Our positive source-restricted schema lacks universal coverage and itself relativizes. There is no identified nonalgebrizing step.

\paragraph{WRAPPER and time-space firewall.}
Fresh completion-time bits affect pre-release prediction, not ordinary post-disclosure correctness of an explicitly supplied SAT instance. No hash hardness is assumed. A bytecode interpreter may have nonlinear branches, DELAY, RECOMP, affine/Tseitin, Horn and 2-SAT routines, but only checked protected writes and one public output. Unchecked COPY, raw STORE or hidden I/O invalidate the theorem. There is no claimed compiler from arbitrary polynomial SAT deciders and no joint time-space contraction; do not combine the verifier's space with another solver's time.

\section{TESTS AND OUTCOME}
Six independently checkpointed Python processes under 29-second individual timeouts, common seed 151015 and final script SHA-256 fc54c0aecdc6a3f7baa7f2f5a442717277bcc9663717806f2201562f432edc00, produced 34,502 exact assertions across 6,393 finite cases, with zero final mismatches. The breakdown was 1,024 exhaustive two-toggle tests (4,416 checks), 1,500 randomized nonlinear paths (19,508), four adversarial examples (518), 3,600 operations in a SPARSE model of a $2^{128}$-cell bank (3,637), all 256 three-variable Boolean truth tables for the circuit-UNSAT equivalence (4,352), and nine explicit full-sweep runs for $2^d$ physically written cells, $2\le d\le10$ (2,071). A symbolic $2^{128}$-cell address space was not materialized. Finite experiments remain COMPUTATIONAL.

\paragraph{Five-track ledger.}
LSPV-1 CRITICALPRISM1: proper clause-deletion answer blindness with source-aware polynomial escape, all-input contraction OPEN. LSPV-2 SEMANTICORBIT1: certified polyhedral symmetry quotient, arbitrary decider-to-proof and full RS5(H) OPEN. LSPV-3 VERTEXCOVERFIBER1: tractable low-exception-cover ternary CSP, common-input UWI non-P obstruction OPEN. LSPV-4 TWOWRITEFRONTIER1: exact fixed-address-word simulation but large time exponent, fast arbitrary graph simulation OPEN. LSPV-5 SYNDROMEGUARD1: exact restricted conservation with exponential physical memory, universal source-grounded certificates and nonalgebrization OPEN.

\paragraph{One next target.}
A charged source-to-certificate transport lemma for UNGUARDED ordinary machine transitions, starting with $x'=x\oplus C(z)$ as a mandatory coNP-hard hostile case. It must not silently introduce runtime guards, assume a short invariant for arbitrary code or claim a P vs NP separation.

\begin{thebibliography}{9}
\bibitem{H} S. Hirahara, Unexpected Hardness Results for Kolmogorov Complexity Under Uniform Reductions, STOC 2020, ECCC TR20-050. \url{https://eccc.weizmann.ac.il/report/2020/050/download}.
\bibitem{B} T. Baker, J. Gill, R. Solovay, Relativizations of the P = ? NP Question, SIAM J. Comput. 4(4), 1975. \url{https://doi.org/10.1137/0204037}.
\bibitem{K} M. Karr, Affine Relationships Among Variables of a Program, Acta Informatica 6, 1976.
\bibitem{M} T. Murata, Petri Nets: Properties, Analysis and Applications, Proc. IEEE 77(4), 1989.
\bibitem{N} G. Necula, Proof-Carrying Code, POPL 1997.
\bibitem{A} S. Aaronson, A. Wigderson, Algebrization: A New Barrier in Complexity Theory, ACM TOCT 2009.
\end{thebibliography}
\end{document}
% FINAL_VERIFICATION_SENTINEL::LSPV5::SYNDROMEGUARD1::2026-10-15_1254::OPEN_P_VS_NP